\section{已登记来源}

\begin{sourcequalitybox}[Appendix use]
附录保留长表用于审计，不承担主文的投资叙事。主文只使用已通过证据等级约束的结论。
\end{sourcequalitybox}

\begin{longtable}{L{1cm}L{2.5cm}L{4.5cm}L{1.8cm}L{1.8cm}L{2.2cm}}
\toprule
\textbf{ID} & \textbf{来源} & \textbf{标题} & \textbf{日期} & \textbf{质量} & \textbf{用途} \\
\midrule
\endfirsthead
\toprule
\textbf{ID} & \textbf{来源} & \textbf{标题} & \textbf{日期} & \textbf{质量} & \textbf{用途} \\
\midrule
\endhead
S01 & 国金证券/慧博 & 再谈PCB的半导体化 & 2026-05-17 & Q2 & 技术主线 \\
S02 & 申万宏源/原创力 & PCB/CCL行业2026策略 & 2026-03 & Q2 & 行业架构 \\
S03 & 国盛证券/发现报告 & 电子2026年度策略 & 2026 & Q2 & 市场空间、M9 \\
S04 & 东吴证券/慧博 & PCB材料端供需缺口 & 2026-05-08 & Q2 & 设备CAPEX \\
S05 & 国海证券/慧博 & 深南电路公司动态 & 2026-05-26 & Q2 & 财务与预测 \\
S06 & 高盛/新浪转载 & 沪电股份管理层调研 & 2026-05-24 & Q2 & 目标价、CAPEX \\
S07 & 摩根大通/新浪转载 & 胜宏科技首次覆盖 & 2026-06-09 & Q2 & H股目标、弹性 \\
S08 & 高盛/新浪转载 & 生益科技AI CCL & 2026-05-22 & Q2 & CCL涨价 \\
S09 & 西部证券/原创力 & 封装基板深度报告 & 2026-05 & Q2 & ABF/BT/玻璃基板 \\
S10 & 东吴证券/迈博 & PCB设备行业总结 & 2026-05-13 & Q2 & 设备耗材 \\
S11 & 申万宏源/中财网 & 华正新材首次覆盖 & 2026-04-03 & Q2 & 财务预测、材料 \\
S12 & 瑞银证券/慧博 & 鹏鼎控股AI PCB新势力 & 2026-05-29 & Q4 & 正文不可用 \\
SEC01 & CMBI/招银国际 & PCB/CCL sector AI-driven upcycle & 2025-11-10 & Q2 & 行业规模、集中度 \\
M01 & 腾讯/新浪快照 & 当前行情快照 & 2026-06-12 & Q2 & 指示性估值 \\
\bottomrule
\caption{来源注册表摘要}
\end{longtable}

\section{高影响主张审计}

\begin{exhibitbox}[附录图表A：主张处理]
\footnotesize
\begin{tabularx}{\textwidth}{X L{3cm}L{3.4cm}}
\toprule
\textbf{主张} & \textbf{分类} & \textbf{报告处理} \\
\midrule
AI服务器升级通过更高层数、低损耗材料、密集互联和背板/载板化内容提升PCB价值。 & broker-stated, cross-source & 作为核心主张，但加入估值纪律。 \\
M7/M8/M9/M10迁移使CCL成为性能瓶颈。 & broker-stated & 支撑材料链核心立场。 \\
胜宏科技H股目标价600港元。 & broker-stated, non-comparable & 列入表格，但不与A股直接比较。 \\
深南电路和华正新材盈利预测。 & broker forecast & 用于指示性PE，不作为AStock预测。 \\
鹏鼎控股AI PCB机会。 & alternative-source verified & HKEX片段和公开PDF支持观察名单，但原UBS全文仍缺，不进入核心结论。 \\
\bottomrule
\end{tabularx}
\end{exhibitbox}

\section{完整报告归档摘要}

\begin{exhibitbox}[附录图表B：核心公司完整报告归档]
\small
\begin{tabularx}{\textwidth}{L{1.8cm}L{2.2cm}X L{2.4cm}}
\toprule
\textbf{公司（代码）} & \textbf{机构} & \textbf{标题} & \textbf{状态} \\
\midrule
沪电股份（002463） & 中邮证券 & 高速交换机业务高增，产能扩张蓄力长期成长 & 已归档并抽取正文 \\
胜宏科技（300476） & 开源证券 & 2026Q1业绩同环比增长，产能建设加快推进中 & 已归档并抽取正文 \\
深南电路（002916） & 太平洋 & AI算力+存储产品共驱25年业绩超预期增长 & 已归档并抽取正文 \\
生益科技（600183） & 太平洋 & 产品结构优化，业绩爆发式增长 & 已归档并抽取正文 \\
华正新材（603186） & 华安证券 & CBF材料阶段性突破，CCL周期修复可期 & 已归档并抽取正文 \\
\bottomrule
\end{tabularx}
\sourcenote{\texttt{workspace/research/semiconductor-pcb-20260612/sources/broker-core-20260615/index.md}}
\end{exhibitbox}

\begin{exhibitbox}[附录图表C：全球投行缺口的替代归档与探测]
\footnotesize
\begin{tabularx}{\textwidth}{L{1.6cm}L{2.0cm}X L{2.3cm}}
\toprule
\textbf{公司（代码）} & \textbf{机构} & \textbf{标题} & \textbf{状态} \\
\midrule
鹏鼎控股（002938） & 浙商证券 & 全球PCB龙头卡位算力赛道，AI驱动第二成长曲线 & 替代报告归档 \\
沪电股份（002463） & 东莞证券 & 技术卡位优势明显，AI算力产品迎收获期 & 替代报告归档 \\
沪电股份（002463） & 海通证券 & 深耕高端PCB市场，AI浪潮迎发展机遇 & 替代报告归档 \\
生益科技（600183） & 西南证券 & 全面拥抱AI浪潮，业绩弹性持续释放 & 替代报告归档 \\
生益科技（600183） & 国金证券 & 高胜率强阿尔法厂商，打破垄断迎成长 & 替代报告归档 \\
沪电股份（002463） & 东吴证券 & 赴港递表加速全球化，谷歌TPU放量迎量价齐升 & 补充探测归档 \\
胜宏科技（300476） & 华金证券 & 聚焦AI服务器高端产品需求，业绩增长动能强劲 & 补充探测归档 \\
行业 & 东吴证券 & AI驱动PCB全面升级：材料、工艺与架构革新引领产业新周期 & 补充探测归档 \\
鹏鼎控股（002938） & 华泰证券 & 卡位AI端侧浪潮，加快算力硬板投入 & 补充探测归档 \\
鹏鼎控股（002938） & 港交所文件 & 业务概览摘录 & 发行文件归档 \\
生益科技（600183） & 招银国际 & 拥抱AI、迈向新增长轨道 & 刷新归档 \\
生益科技（600183） & 招银国际 & PCB表现优于CCL，利润率环比改善 & 刷新归档 \\
胜宏科技（300476） & 摩根大通/Reportify & AI热潮驱动的PCB价值量增长周期主要受益者 & 可见转写归档 \\
胜宏科技（300476） & 摩根大通/新浪转载 & AI引爆PCB价值量提升，给予超配评级 & 转载正文归档 \\
生益科技（600183） & 高盛/新浪转载 & 高端CCL覆铜板价值量爆发，目标价大幅上调 & 转载正文归档 \\
生益科技（600183） & 花旗/新浪转载 & 覆铜板涨价潮远未结束，目标价被大幅上调 & 转载正文归档 \\
鹏鼎控股（002938） & 天风证券 & 加大AI PCB投入，软硬板发力掌握AI云网端成长机遇 & 刷新归档 \\
鹏鼎控股（002938） & 开源证券 & 2025H1业绩超预期，AI云-管-端PCB全链条布局 & 刷新归档 \\
鹏鼎控股（002938） & 民生证券 & 打造AI终端的创新底座 & 30页深度报告归档 \\
鹏鼎控股（002938） & 兴业证券 & 端侧AI渐行渐近，PCB龙头受益软硬板量价双升趋势 & 历史报告归档 \\
沪电股份（002463） & 信达证券 & 2025半年报点评：AI服务器和交换机延续高景气 & 刷新归档 \\
沪电股份（002463） & 天风证券 & 25年逐季度延续高增态势，AI浪潮带动产品需求火热 & 刷新归档 \\
华正新材（603186） & 浙商证券 & 行业结构性复苏确立拐点，产品结构高端化定调成长 & 刷新归档 \\
华正新材（603186） & 浙商证券 & CBF膜国产化主力军 & 历史报告归档 \\
华正新材（603186） & 海通国际 & 算力侧材料国产替代进行时 & 历史报告归档 \\
\bottomrule
\end{tabularx}
\sourcenote{\texttt{data/global\_broker\_gap\_fallback\_summary.md}; \texttt{data/global\_broker\_original\_pdf\_probe\_20260616.md}; \texttt{data/original\_pdf\_refresh\_20260616.md}; \texttt{data/global\_broker\_current\_recheck\_20260617.md}; these visible repost/transcript files do not replace missing UBS/JPM/Goldman/Citi original PDFs}
\end{exhibitbox}

\section{具名客户传闻隔离}

\begin{sourcequalitybox}[Unverified named-customer claims]
市场上存在关于 NVIDIA、Google TPU/ASIC、Rubin、Apple、M9/M10认证等具名客户或平台份额的说法。除非来自公司公告、官方IR逐字稿、原始券商PDF或客户/供应商披露，本报告不将其作为确认事实。Apple 官方 supply-chain 页面和旧 supplier-list PDF 路径已探测，但本轮未取得可审计的鹏鼎/臻鼎供应商条目或收入拆分。2026年6月17日复查仍只找到财富号、股吧、雪球、头条、证券星、微博等二级/社媒/转载来源的具体份额和订单数字，不能升格为确认收入拆分。
\sourcenote{\texttt{data/named\_customer\_rumor\_registry.md}; \texttt{data/apple\_customer\_side\_official\_probe\_20260616.md}; \texttt{data/current\_public\_source\_recheck\_20260617.md}}
\end{sourcequalitybox}

\section{具名平台收入拆分确认包}

\begin{exhibitbox}[附录图表D：IR / paid-data confirmation fields]
\scriptsize
\begin{tabularx}{\textwidth}{L{1.5cm}L{1.7cm}X X}
\toprule
\textbf{代码} & \textbf{公司} & \textbf{需确认字段} & \textbf{公开边界} \\
\midrule
002463 & 沪电股份 & NVIDIA/Rubin/GB系列、Google TPU/ASIC、AMD/海外云、国产算力、高速交换机/光模块平台收入；相关ASP、订单和毛利。 & 已有AI服务器/HPC和高速交换机收入基线；互动易对具体厂商和客户合作以商业政策限制为由不披露。 \\
300476 & 胜宏科技 & NVIDIA/Blackwell/Rubin/GB200/GB300、Google TPU/ASIC、Microsoft/Amazon ASIC、UBB/OAM/midplane/正交背板、1.6T光模块mSAP收入。 & 已有ASIC进展、mSAP订单饱满和高层/HDI能力证据；具体客户名称和业务细节仍受商业政策限制。 \\
002916 & 深南电路 & AI加速卡PCB、高速交换机、光模块PCB、服务器/存储PCB、FC-BGA/ABF/BT载板收入按客户/平台拆分。 & 已有PCB/载板分部和AI硬件需求证据；具体企业合作问题按商业保密原则不回复。 \\
600183 & 生益科技 & M8/M9/M10收入，GPU/ASIC AI服务器材料、国产算力、交换机/光模块材料收入，具名客户认证和批量供货状态。 & 已有高频高速材料和国内外终端GPU/AI项目合作证据；不披露具体客户或M8/M9/M10收入占比。 \\
603186 & 华正新材 & 高速CCL按AI服务器/交换机/光模块/天线拆分，BT/类BT和CBF/CBF-RCC按客户/应用拆分，国产算力验证和爬坡收入。 & 已有AI服务器/交换机/光模块材料方向和CBF/BT进展；不披露具名客户、ASP、出货或订单金额。 \\
\bottomrule
\end{tabularx}
\sourcenote{\texttt{data/named\_platform\_customer\_revenue\_split\_confirmation\_pack\_20260618.md}; \texttt{data/named\_platform\_customer\_revenue\_split\_confirmation\_pack\_20260618.json}}
\end{exhibitbox}

\begin{sourcequalitybox}[Completion test for this blocker]
具名平台收入拆分的完成条件不是新增供应商关系证据，而是取得可归因的客户/平台收入拆分，或取得公司、客户、供应商明确确认“该颗粒度不可披露”。确认包已经列出五家公司联系人、待确认字段、公开收入基线和披露边界；在没有直接确认或付费数据库前，本报告仍将该项保持为外部数据需求。
\end{sourcequalitybox}
